%% ===========================================================================
%%  Abstract + index terms (IEEE format).
%%  The ACM CCS concept block (\begin{CCSXML} ... \ccsdesc) has been removed:
%%  IEEE journals do not use CCS; keywords become \begin{IEEEkeywords}.
%%
%%  NOTE: this is still the ACM-submission abstract.  See
%%  docs/reports/TCSVT_migration_notes.md §2 for the recommended rewrite that
%%  matches the TCSVT positioning (systems/analysis framing, controllable
%%  preservation--edit trade-off, new evidence).
%% ===========================================================================

\begin{abstract}
Text-guided image editing aims to modify real images according to natural language instructions while preserving non-target visual content. Existing diffusion- and rectified flow-based methods face a fundamental trade-off between editing effectiveness and source fidelity, stemming from inversion errors inherent to continuous latent trajectories and cross-modal feature entanglement that causes unintended global changes. We present a \textbf{training-free text-guided image editing framework based on Visual Autoregressive Modeling (VAR)}, which sidesteps both challenges by operating on discrete, scale-wise token generation rather than continuous noise removal. Since no trajectory exists to invert, tokens can be selectively replaced at the sampling stage without altering the global generation process. Our method leverages cross-modal attention maps between text and visual tokens to automatically derive high-attention focus masks and low-attention preserve masks at each generation scale, eliminating the need for external segmentation tools or per-category hyperparameter tuning. Binarizing those maps is the step that governs the preservation--edit trade-off, and we show that rank-based and mean-relative thresholds both fail for the same structural reason; we instead threshold an absolute energy level on the normalized attention map with an instance-adaptive correction, whose parameters are calibrated so that the realized threshold distribution is centred on a chosen operating point rather than tuned by inspection. For real image editing, the source image is encoded into discrete per-scale bitwise tokens via BSQ-VAE, enabling direct token-level background substitution that faithfully preserves source content without inversion or continuous feature injection. Beyond the method, we systematically isolate which components govern the trade-off: the base threshold traces a smooth monotone frontier and is the single control worth exposing, the instance-adaptive correction is a small refinement on top of it rather than a second lever, the target-stream pass contributes structurally rather than parametrically, and one apparent hyperparameter provably cannot affect the output. On PIE-Bench our method attains the best SSIM and the best CLIP whole-image similarity among diffusion-, flow- and autoregressive-based training-free baselines, is second on LPIPS and on edited-region CLIP, and remains in the leading group on human-preference scores while using a 2B-parameter backbone --- six times smaller than the flow-based baselines --- in an inversion-free pipeline that runs in about 2.5 seconds per image.
\end{abstract}

\begin{IEEEkeywords}
Image editing, text-guided image manipulation, visual autoregressive modeling, training-free editing, inversion-free editing, cross-modal attention, discrete token representation.
\end{IEEEkeywords}
